% GENERATED FILE. Edit canonical lessons, then run npm run generate:books.
% canonical-source-hash: fnv1a64:a3a7886853a41c97
% canonical-lessons: HI-C494-timardari, HI-C494-god, HI-C494-litana, HI-C494-sekna, HI-C494-orhana

\chapter{Looking After Someone}
\label{ch:hi-looking-after-someone}
\hlchaptermodality{494}

\begin{chapteropening}
\textbf{By the end of this chapter:} \emph{I can say how I look after someone who is ill.}

Complete the last lesson of chapter 494: I can say how I look after someone who is ill.
\end{chapteropening}

\section[\texorpdfstring{\hi{तीमारदारी} (tīmārdārī)}{तीमारदारी (tīmārdārī)}]{\hi{तीमारदारी} (tīmārdārī) — nursing the sick}

\label{lesson:HI-C494-timardari}



\begin{quote}
Before the new one: say the Hindi for cancer, then the Hindi for to stiffen.
\end{quote}

\subsection*{You'll want to know: \hi{तीमारदारी}}
\textbf{\hi{तीमारदारी}} — \emph{tīmārdārī} — \textquotedblleft{}nursing the sick\textquotedblright{}.

\begin{quote}
\hi{बीमारी} \hi{में} \hi{मीरा} \hi{ने} \hi{अपनी} \hi{माँ} \hi{की} \hi{बहुत} \hi{तीमारदारी} \hi{की।}
\end{quote}

\begin{quote}
\emph{bīmārī meṁ Mīrā ne apnī māṁ kī bahut tīmārdārī kī.}
\end{quote}

\begin{quote}
Mira nursed her mother devotedly through her illness.
\end{quote}

Feminine.

\begin{cousinweb}
\textbf{\hi{तीमारदारी}} (\emph{tīmārdārī}) is usually traced to Persian \emph{tīmārdārī}, \textquotedblleft{}tending the sick\textquotedblright{}, built on \emph{tīmār}, \textquotedblleft{}care of the sick\textquotedblright{}, and \emph{-dārī}, \textquotedblleft{}keeping\textquotedblright{}. \emph{tīmārdār} is the person at the bedside.
\end{cousinweb}

\begin{grammarlens}[title={What you've built}]
One more word for describing what you lived and felt.
\end{grammarlens}

\subsection*{Your turn}
\begin{itemize}
\raggedright
  \item \emph{Say it:} \emph{tīmārdārī}
  \item \emph{Say it:} \emph{tīmārdārī}, once more
  \item \emph{Say it:} \emph{bīmārī meṁ Mīrā ne apnī māṁ kī bahut tīmārdārī kī.} — Mira nursed her mother devotedly through her illness.
  \item \emph{Recall:} say the Hindi for cancer, then the Hindi for to stiffen, then say \emph{tīmārdārī} again
\end{itemize}

\begin{tcolorbox}[breakable,colback=teal!4,colframe=teal!35!black,title={Before you move on}]
What does \hi{तीमारदारी} mean? (\textquotedblleft{}Nursing the sick\textquotedblright{}.) Say it once more.
\end{tcolorbox}
\section[\texorpdfstring{\hi{गोद} (god)}{गोद (god)}]{\hi{गोद} (god) — a lap}

\label{lesson:HI-C494-god}



\begin{quote}
Before the new one: say the Hindi for to stiffen, then the Hindi for nursing the sick.
\end{quote}

\subsection*{You'll want to know: \hi{गोद}}
\textbf{\hi{गोद}} — \emph{god} — \textquotedblleft{}a lap\textquotedblright{}.

\begin{quote}
\hi{बच्चा} \hi{माँ} \hi{की} \hi{गोद} \hi{में} \hi{सो} \hi{गया।}
\end{quote}

\begin{quote}
\emph{baccā māṁ kī god meṁ so gayā.}
\end{quote}

\begin{quote}
The baby fell asleep in its mother's lap.
\end{quote}

Feminine.

\begin{cousinweb}
\textbf{\hi{गोद}} (\emph{god}) is a tadbhava: CDIAL lists Hindi \emph{god}, \textquotedblleft{}lap\textquotedblright{}, under \emph{*gōddī}, a form it reconstructs (CDIAL 4284). \emph{god lenā}, \textquotedblleft{}to take into one's lap\textquotedblright{}, is to adopt a child.
\end{cousinweb}

\begin{grammarlens}[title={What you've built}]
One more word for describing what you lived and felt.
\end{grammarlens}

\subsection*{Your turn}
\begin{itemize}
\raggedright
  \item \emph{Say it:} \emph{god}
  \item \emph{Say it:} \emph{god}, once more
  \item \emph{Say it:} \emph{baccā māṁ kī god meṁ so gayā.} — The baby fell asleep in its mother's lap.
  \item \emph{Recall:} say the Hindi for to stiffen, then the Hindi for nursing the sick, then say \emph{god} again
\end{itemize}

\begin{tcolorbox}[breakable,colback=teal!4,colframe=teal!35!black,title={Before you move on}]
What does \hi{गोद} mean? (\textquotedblleft{}A lap\textquotedblright{}.) Say it once more.
\end{tcolorbox}
\section[\texorpdfstring{\hi{लिटाना} (liṭānā)}{लिटाना (liṭānā)}]{\hi{लिटाना} (liṭānā) — to lay someone down}

\label{lesson:HI-C494-litana}



\begin{quote}
Before the new one: say the Hindi for nursing the sick, then the Hindi for a lap.
\end{quote}

\subsection*{You'll want to know: \hi{लिटाना}}
\textbf{\hi{लिटाना}} — \emph{liṭānā} — \textquotedblleft{}to lay someone down\textquotedblright{}.

\begin{quote}
\hi{नर्स} \hi{ने} \hi{मरीज़} \hi{को} \hi{बिस्तर} \hi{पर} \hi{लिटाया।}
\end{quote}

\begin{quote}
\emph{nars ne marīz ko bistar par liṭāyā.}
\end{quote}

\begin{quote}
The nurse laid the patient down on the bed.
\end{quote}

Past, with \hi{ने} (he or she …): \hi{उसने} \hi{बच्चे} \hi{को} \hi{लिटाया}, \emph{usne bacce ko liṭāyā}.

\begin{cousinweb}
\textbf{\hi{लिटाना}} (\emph{liṭānā}) is the causative partner of \emph{leṭnā}, \textquotedblleft{}to lie down\textquotedblright{}, which CDIAL lists under Sanskrit \emph{lēṭyáti} (CDIAL 11109); the vowel shortens before the causative ending.
\end{cousinweb}

\begin{grammarlens}[title={What you've built}]
One more word for describing what you lived and felt.
\end{grammarlens}

\subsection*{Your turn}
\begin{itemize}
\raggedright
  \item \emph{Say it:} \emph{liṭānā}
  \item \emph{Say it:} \emph{liṭānā}, once more
  \item \emph{Say it:} \emph{nars ne marīz ko bistar par liṭāyā.} — The nurse laid the patient down on the bed.
  \item \emph{Recall:} say the Hindi for nursing the sick, then the Hindi for a lap, then say \emph{liṭānā} again
\end{itemize}

\begin{tcolorbox}[breakable,colback=teal!4,colframe=teal!35!black,title={Before you move on}]
What does \hi{लिटाना} mean? (\textquotedblleft{}To lay someone down\textquotedblright{}.) Say it once more.
\end{tcolorbox}
\section[\texorpdfstring{\hi{सेंकना} (sẽknā)}{सेंकना (sẽknā)}]{\hi{सेंकना} (sẽknā) — to warm (a sore part); to toast}

\label{lesson:HI-C494-sekna}



\begin{quote}
Before the new one: say the Hindi for a lap, then the Hindi for to lay someone down.
\end{quote}

\subsection*{You'll want to know: \hi{सेंकना}}
\textbf{\hi{सेंकना}} — \emph{sẽknā} — \textquotedblleft{}to warm (a sore part); to toast\textquotedblright{}.

\begin{quote}
\hi{माँ} \hi{ने} \hi{गरम} \hi{पानी} \hi{की} \hi{बोतल} \hi{से} \hi{मेरा} \hi{पेट} \hi{सेंका।}
\end{quote}

\begin{quote}
\emph{māṁ ne garam pānī kī botal se merā peṭ sẽkā.}
\end{quote}

\begin{quote}
Mother warmed my stomach with a hot-water bottle.
\end{quote}

Past, with \hi{ने} (he or she …): \hi{उसने} \hi{घुटना} \hi{सेंका}, \emph{usne ghuṭnā sẽkā}.

\begin{cousinweb}
\textbf{\hi{सेंकना}} (\emph{sẽknā}) is a tadbhava: CDIAL lists Hindi \emph{sẽknā}, \textquotedblleft{}to warm, foment, toast\textquotedblright{}, under \emph{*sēkk}, a form it reconstructs (CDIAL 13581). Bread is toasted and a sore knee warmed with the same verb.
\end{cousinweb}

\begin{grammarlens}[title={What you've built}]
One more word for describing what you lived and felt.
\end{grammarlens}

\subsection*{Your turn}
\begin{itemize}
\raggedright
  \item \emph{Say it:} \emph{sẽknā}
  \item \emph{Say it:} \emph{sẽknā}, once more
  \item \emph{Say it:} \emph{māṁ ne garam pānī kī botal se merā peṭ sẽkā.} — Mother warmed my stomach with a hot-water bottle.
  \item \emph{Recall:} say the Hindi for a lap, then the Hindi for to lay someone down, then say \emph{sẽknā} again
\end{itemize}

\begin{tcolorbox}[breakable,colback=teal!4,colframe=teal!35!black,title={Before you move on}]
What does \hi{सेंकना} mean? (\textquotedblleft{}To warm (a sore part); to toast\textquotedblright{}.) Say it once more.
\end{tcolorbox}
\section[\texorpdfstring{\hi{ओढ़ाना} (oṛhānā)}{ओढ़ाना (oṛhānā)}]{\hi{ओढ़ाना} (oṛhānā) — to cover someone (with a blanket)}

\label{lesson:HI-C494-orhana}



\begin{quote}
Before the new one: say the Hindi for to lay someone down, then the Hindi for to warm (a sore part).
\end{quote}

\subsection*{You'll want to know: \hi{ओढ़ाना}}
\textbf{\hi{ओढ़ाना}} — \emph{oṛhānā} — \textquotedblleft{}to cover someone (with a blanket)\textquotedblright{}.

\begin{quote}
\hi{माँ} \hi{ने} \hi{बच्चे} \hi{को} \hi{कंबल} \hi{ओढ़ाया।}
\end{quote}

\begin{quote}
\emph{māṁ ne bacce ko kambal oṛhāyā.}
\end{quote}

\begin{quote}
Mother tucked a blanket round the child.
\end{quote}

Past, with \hi{ने} (he or she …): \hi{उसने} \hi{कंबल} \hi{ओढ़ाया}, \emph{usne kambal oṛhāyā}.

\begin{cousinweb}
\textbf{\hi{ओढ़ाना}} (\emph{oṛhānā}) is the causative partner of \emph{oṛhnā}, \textquotedblleft{}to cover oneself\textquotedblright{}, which CDIAL lists under \emph{*ōḍḍh}, a form it reconstructs (CDIAL 2547): to wrap someone else up.
\end{cousinweb}

\begin{grammarlens}[title={What you've built}]
One more word for describing what you lived and felt.
\end{grammarlens}

\subsection*{Your turn}
\begin{itemize}
\raggedright
  \item \emph{Say it:} \emph{oṛhānā}
  \item \emph{Say it:} \emph{oṛhānā}, once more
  \item \emph{Say it:} \emph{māṁ ne bacce ko kambal oṛhāyā.} — Mother tucked a blanket round the child.
  \item \emph{Recall:} say the Hindi for to lay someone down, then the Hindi for to warm (a sore part), then say \emph{oṛhānā} again
\end{itemize}

\begin{tcolorbox}[breakable,colback=teal!4,colframe=teal!35!black,title={Before you move on}]
What does \hi{ओढ़ाना} mean? (\textquotedblleft{}To cover someone (with a blanket)\textquotedblright{}.) Say it once more.
\end{tcolorbox}
